\documentclass[10pt,a4paper]{amsart}
\usepackage[margin=19mm]{geometry}
\usepackage{amsmath,amssymb,mathtools}
\usepackage[scr=rsfs,cal=euler]{mathalfa}
\usepackage{xfrac}
\usepackage[unicode,hidelinks,bookmarks=false]{hyperref}
\newcommand{\ie}{i.e.\ }
\newcommand{\cf}{cf.\ }
\newcommand{\resp}{respectively\ }
\newcommand{\Subsection}{\subsection}
\providecommand{\nobreakdash}{\nobreak\hbox{-}}
\setlength{\emergencystretch}{2em}
\multlinegap0pt
\usepackage{amssymb}
\usepackage{mathtools}
\usepackage{dsfont}
\usepackage[shortlabels]{enumitem}
\newtheorem{proposition}{Proposition}
\def\de{\delta}
\title{Schur's theorem and its relation to the closure properties of the non-abelian tensor product}
\author{Guram Donadze, Manuel Ladra, Pilar P\'aez-Guill\'an}
\date{Source: arXiv:2601.19764v1 (CC BY 4.0)}
\begin{document}
\maketitle

\noindent\emph{Source and licence.} Schur's theorem and its relation to the closure properties of the non-abelian tensor product. Version arXiv:2601.19764v1 (CC BY 4.0), distributed by arXiv under the Creative Commons Attribution 4.0 International licence, http://creativecommons.org/licenses/by/4.0/, as recorded in the arXiv licence metadata for this exact version and shown on the abstract page https://arxiv.org/abs/2601.19764v1. Full licence notice: \texttt{licenses/arxiv-2601.19764-CC-BY-4.0.txt}.\par\smallskip

\noindent\textit{Editorial scope.} Counted unit: the article's proposition prop3 (source0.tex lines 523--526). The article's complete original proof of it is delivered with it (source0.tex lines 527--550, one \texttt{proof} environment of 24 lines). No mathematics is written, repaired or re-derived: the statement and the proof are the article's own lines, in the article's own notation, and no retained line is substituted or re-flowed.  Retained uncounted as the source-local context the proof needs: lines 248--251.  Nothing was omitted from any retained span, so the package carries no omission record.  No retained line contains a citation, so the wrapper carries no bibliography and no bibliography entry.  No retained line contains a figure environment, an image-inclusion command or drawing code, so no image is delivered and the package declares no asset dependency.  Editorial formatting only, in addition: an AMS \texttt{amsart} preamble that restates the article's own declarations and macros used by the retained lines, namely the environments \texttt{proposition} on separate counters, together with the standard packages those lines need.  The retained environments print in this document as Proposition 1 in order of appearance, whereas the article prints the same items with the numbering of its own full document.  The complete native source of the article as distributed in the arXiv e-print for this exact version is bundled unchanged as \texttt{sources/193483/source0.tex}.
\begin{proposition}\label{paper5}
Let $G$ and $H$ be groups acting on each other compatibly. If $G$ and $H$ are finitely generated, then $G\otimes H$ is also finitely generated if and only if so are
$D_H(G)$ and $D_G(H)$.
\end{proposition}

\begin{proposition} \label{prop3}
Let $H$ be a finitely generated group and $1\to N \to G \to H \to 1$ be a central extension of $H$. Fix an integer $n\geq 1$ and suppose that $\gamma_n(G)$ is a finitely
generated group. Then, $\gamma_{n+1}(G)$ is finitely generated if and only if $\gamma_{n+1}(H)$ is finitely generated.
\end{proposition}

\begin{proof} The idea of this demonstration is to show that if  $\gamma_{n+1}(H)$ is finitely generated, then $\gamma_{n+1}(G)$ is finitely generated too. We have the following extensions
of groups:
\[
1\to N \to G \to H \to 1 ,
\]
and
\[
1\to N_n \to \gamma_n(G) \to \gamma_n(H) \to 1,
\]
where $N_n = N\cap \gamma_n(G)$. Therefore, considering the non-abelian tensor products we get the following exact sequence:
\[
N\otimes \gamma_n(G) \times N_n\otimes G \to G\otimes \gamma_n(G) \to H\otimes \gamma_n(H) \to 1 .
\]
This sequence leads to the isomorphism:
\[
\frac{ G\otimes \gamma_n(G)}{X} \cong  H\otimes \gamma_n(H),
\]
where $X$ is the image of $N\otimes \gamma_n(G) \times N_n\otimes G$ into $G\otimes \gamma_n(G)$. Now, let
$\de \colon G\otimes \gamma_n(G)\to \gamma_{n+1}(G)$ be the homomorphism given by $g\otimes g' \mapsto [g, g']$ for
each $g\in G$ and $g' \in \gamma_{n}(G)$. Since $N$ and $N_n$ are contained in the center of $G$, it holds that $\de (X)=1$. Therefore, $\de$ induces an epimorphism
from $ H\otimes \gamma_n(H)$ to $\gamma_{n+1}(G)$, and it suffices to show that $H\otimes \gamma_n (H)$ is finitely generated. Note that
$D_{H}(\gamma_n(H))=\gamma_{n+1}(H)$ and $D_{\gamma_n (H)}(H)=\gamma_{n+1}(H)$. By Proposition~\ref{paper5}, $H\otimes \gamma_n (H)$ is
finitely generated. Hence, $\gamma_{n+1}(G)$ is finitely generated too.
\end{proof}

\end{document}
